\section{The source carrying section and its exact cost}
\label{sec:ftb-carrying-section}
This section follows the construction in
\texttt{ExistsCarryingSectionOrders} and
\texttt{CarryingSectionIntegralFrame}.  It keeps the same component,
level and reference weight throughout.  The analytic hypotheses used below
are the actual estimates for that component; no estimate for a different
cover or for a larger space is substituted.

Let \(E\) be the CM host, \(f=|\Psi|\), \(r=p+q\geq1\), and let \(K\)
be the coefficient field, \(d_K=[K:\mathbf Q]\).
Fix the good uniformizations \(U_\sigma\), their frame dictionaries, a
component \(Z\), and the source metric parameter \(x\).
Write \(H\) for the source Hodge weight, with rows
\(H_{\tau i}=0\) for \(i<p\) and \(H_{\tau i}=-1\) for \(i\geq p\).
Thus
\[
 u'=u-3rH,\qquad
 G=\operatorname{wt}(u,0),\quad
 G_0=\operatorname{wt}(8(u-u'),0)=\operatorname{wt}(24rH,0).
\]
All section weights in this section are integer weights.
The scalar exponent is even wherever the source norm multiplication
law or the integer-reading law is used.

For a finite-place function \(e=(e_w)\) of finite support and a
conjugation-invariant function \(t=(t_\sigma)\), put
\[
 \operatorname{cost}(e,t)=\frac1{d_K}
  \left(\sum_w e_w\log Nw+\sum_\sigma t_\sigma\right).
\]
For a recurrent valuation \(v\) over \(w\), let
\[
 \operatorname{depth}_Z(G,v)
   =\inf\{v(s):s\in\operatorname{latOn}(U,Z,G)\}.
\]
Here \(\operatorname{latOn}\) is the source lattice of integral cusp
sections supported on \(Z\), not a lattice chosen after \(v\).
A section \(\theta\) \emph{carries \(u'\) to \(u\)} if it has weight
\(G_0\), is nonzero at the generic point of \(Z\), and, for every
\(k\geq0\),
\[
 \theta^k\operatorname{cuspOn}(U,Z,\operatorname{wt}(8ku',0))
 \subseteq
 \operatorname{cuspOn}(U,Z,\operatorname{wt}(8ku,0)).
 \tag{C1}\label{eq:ftb-carries}
\]
The finite inequality required later is
\[
 8\operatorname{depth}_Z(G,v)
       \leq v(\theta)+8e_w                         \tag{C2}\label{eq:ftb-finbd}
\]
whenever some integral section of weight \(G\) has finite \(v\).
The infinite inequality is
\[
 \operatorname{homN}_{G_0}(\theta;y)^2
   \leq e^{16t_\sigma}\,
           \operatorname{Berg}_{Z,G}(y)^8
 \quad (y\text{ a frame on }Z).                    \tag{C3}\label{eq:ftb-infbd}
\]
The Bergman function is that of the specified reference space and frame
dictionary.  Statements with value \(+\infty\) have the evident
extended-real interpretation.

\subsection{The finite theta inequality}
Here is the elementary lattice calculation used in the finite
theta estimate.  Let \(Y\) be a finitely generated
\(\mathbf Z\)-module with an \(\mathcal O_E\)-action.  Conjugation
preserves \(\mathcal O_E\); choose \(w_0\in\mathcal O_E\) with
\(w_0+\bar w_0=1\).  Let \(\mathscr B\) be the real vector space of
symmetric \(\mathbf Z\)-bilinear forms satisfying
\(\gamma(ca,b)=\gamma(a,\bar c b)\).
For positive semidefinite \(\beta\in\mathscr B\), define
\[
 \phi_C(\beta)=
 \inf_{\ell}\ell(\beta),                           \tag{C4}\label{eq:ftb-phic}
\]
where \(\ell:\mathscr B\to\mathbf R\) is linear, is integer-valued
on integer-valued forms, and is strictly positive on every nonzero
positive semidefinite form.  This is the source definition of
\(\phi_C\); it is not the minimum of a quadratic form.

\begin{lemma}[Spanning and coset minima]\label{lem:ftb-coset-minimum}
Let \(Q:Y\to\mathbf Z_{\geq0}\) satisfy
\(\beta(a,a)=2Q(a)\).  Suppose \(m\geq2\), and
\(y_1,\ldots,y_n\) generate \(Y/m^2Y\) as an
\(\mathcal O_E\)-module.  If \(b_i\in\mathbf Q\cup\{+\infty\}\)
satisfy, for every \(q\in\mathbf Q\),
\[
 \bigl[q\leq Q(mz+y_i)/m^2\text{ for every }z\in Y\bigr]
       \ \Longrightarrow\ q\leq b_i,
\]
then
\[
                 \phi_C(\beta)/m^2\leq\sum_i b_i .
 \tag{C5}\label{eq:ftb-coset-minimum}
\]
\end{lemma}
\begin{proof}
Choose \(z_i\) minimizing the nonnegative integer
\(Q(y_i+mz)\).  Let \(X\) be the \(\mathcal O_E\)-span of
\(y_i+mz_i\).  The generation assumption gives \(Y=X+mY\).
The finitely generated abelian quotient \(A=Y/X\) satisfies \(A=mA\).
The determinant trick applied to generators of \(A\) supplies a
nonzero integer \(s\equiv1\pmod m\) with \(sA=0\):
write each generator as \(m\) times a combination of the generators
and take the determinant of \(I-mC\).  Its determinant is
\(1\bmod m\), hence is nonzero.  Thus every element of \(Y\) has a
positive integer multiple in \(X\).

For \(x_i=y_i+mz_i\) define
\(\ell(\gamma)=\sum_i\gamma(w_0x_i,x_i)\).
Symmetry, balance, and \(w_0+\bar w_0=1\) give
\(2\gamma(w_0a,a)=\gamma(a,a)\).
The functional \(\ell\) is integer-valued on integral forms.
It is positive on every nonzero positive semidefinite form:
if its nonnegative summands all vanish, then
\(\gamma(x_i,x_i)=0\).  Positivity of
\(\gamma(nx_i+y,nx_i+y)\) for all integers \(n\) forces
\(\gamma(x_i,y)=0\).  Balance extends this vanishing to \(X\), and
the positive-multiple property extends it to \(Y\).  This would give
\(\gamma=0\).  Thus \(\ell\) occurs in \eqref{eq:ftb-phic}, and
\[
 0\leq\phi_C(\beta)
 \leq\sum_i\beta(w_0x_i,x_i)=\sum_iQ(x_i).
\]
The infimum is bounded below by zero because \(\beta\) is
positive semidefinite.  The hypothesis at the rational number
\(Q(x_i)/m^2\), followed by addition, proves \eqref{eq:ftb-coset-minimum}.
\end{proof}

The geometric input for applying this lemma is precise.  On the actual
complete-DVR PEL uniformization of an algebraic integral frame, take
the primitive integral theta section, its expansion of nonnegative
\(\nu\), the cubically normalized origin value, and the odd-level
Heisenberg lifts compatible with the full level structure.
For \(m\) odd the expansion formula and its valuation convergence law
give
\[
 \left[\forall\chi,\ q\leq
      Q(m\chi+y_a)/m^2\right]
       \Longrightarrow q\leq v(\vartheta_a) .
 \tag{C6}\label{eq:ftb-theta-coset-val}
\]
Indeed, the normalized theta value is the convergent sum of its
expansion terms in that coset.  Nonnegative \(\nu\) bounds every
term below by \(q\); the nonarchimedean sum law bounds the sum below
by \(q\).  The lifts of a tuple spanning the full level module
give the \(y_a\)'s spanning \(Y/m^2Y\).
These are statements about the same integral theta section and its
normalized values.  Integral expansions or torsion lifts on a
different model would not suffice.

Apply \eqref{eq:ftb-theta-coset-val} and
Lemma~\ref{lem:ftb-coset-minimum} with \(m=R^2\).
An integral linear combination of products of the corresponding
theta constants has valuation at least \(\phi_C/R^4\).
Each product has this lower bound, its integer coefficient has
nonnegative valuation, and the valuation of a finite sum is at least
the minimum of those bounds.  Finite extension and comparison of the
two semiabelian models preserve the normalized value.  Consequently
the source theta rule \(f\), for any tuple spanning the level
\(R^4\) module, satisfies
\[
             v(f)\geq\phi/R^4                       \tag{C7}\label{eq:ftb-tau-bound}
\]
at a point carrying the actual source datum \(\operatorname{HasPhi}(\phi)\).
Here \(\phi\geq0\); if \(f\) vanishes, the assertion holds with
value \(+\infty\).
This proves the finite calculation from the explicitly stated
same-model expansion and normalized-value inputs.  Their geometric
construction remains the integral-uniformization input, not an
additional claim of this calculation.

\subsection{The section \(f^{24}KS\)}
Assume \(5^4\mid M\), and supply the source admissible level-\(5^4\)
tuple and theta rule \(f\) nonzero on \(Z\).  This nonvanishing input
is the actual shifted-null-value construction of
\texttt{exists\_isGenRule\_nonzeroOn}; its scope is retained below.
The admissible tuple generates the full level module.
Let
\[
 a_\sigma(y)=
       \operatorname{homN}_{\operatorname{wtZ}(0,2r)}(f^2;y).
\]
Retain the following two estimates, uniformly for every such generating
tuple and its theta rule, on the same component and cover:
\[
 a_\sigma(y)\leq A,\qquad
 a_\sigma(y)^3\leq e^\beta\operatorname{Berg}_{Z,G}(y),
 \quad \beta\geq0.                                 \tag{C8}\label{eq:ftb-carry-analytic}
\]
Their source producers are listed below.

Let \(K_0\) be the \(3r\)-th power of the weight-four key unit,
viewed as a twisting section.  Its weight is
\(\operatorname{wt}(0,-12r)\); its finite lower bounds are
nonnegative, and its constant norm factors satisfy
\[
 \log c_{K_0,\sigma}\leq
 \kappa:=3r\,(2fr\log(2\pi)).
\]
Let \(S\) be the flat constant section of weight
\(\operatorname{wt}(\operatorname{par}(-12r),0)\), with finite
values \(\rho_w\), positive constant norm factors \(c_{S,\sigma}\),
and
\[
 \deg S
 =12r\,f\,\frac{p-q}{2}\omega
 =\frac1{d_K}\left(\sum_w\rho_w\log Nw
                         -\sum_\sigma\log c_{S,\sigma}\right).
 \tag{C9}\label{eq:ftb-flat-carry-deg}
\]
Both twisting sections have inverses on \(Z\).

\begin{proposition}[The explicit carrying section]
\label{prop:ftb-carry-explicit}
Under these hypotheses put
\[
 \theta=f^{24}K_0S,\qquad
 e_w=-\rho_w/8,\qquad
 t_\sigma=\beta/2+\kappa/8+\log c_{S,\sigma}/8.
 \tag{C10}\label{eq:ftb-carry-choice}
\]
Then \(\theta\) carries \(u'\) to \(u\).  For every actual
algebraic integral frame with source
\(\operatorname{HasPhi}(\phi)\), over \(w\) with ramification
normalizer \(e_0>0\),
\[
 \frac8{2500}\frac{\phi}{e_0}-8e_w
           \leq v_{\mathrm{norm}}(\theta).          \tag{C11}\label{eq:ftb-carry-finite}
\]
It satisfies \eqref{eq:ftb-infbd}, and its cost is exactly
\[
          \operatorname{cost}(e,t)
                =\beta/2+\kappa/8-\deg S/8.         \tag{C12}\label{eq:ftb-carry-cost}
\]
\end{proposition}
\begin{proof}
The source identity
\[
 12\,\operatorname{wtZ}(0,2r)
  +\operatorname{wt}(0,-12r)
  +\operatorname{wt}(\operatorname{par}(-12r),0)
       =\operatorname{wt}(24rH,0)
\]
gives the required weight.  The generic value is nonzero because
\(f\) is nonzero and \(K_0,S\) are invertible on \(Z\).
The rows of this weight are antitone.
The integer-reading/bounded-section multiplication theorem therefore
sends any even-weight cusp section supported on \(Z\) to a cusp
section after multiplication by \(\theta\): its norm is bounded by
\(A^{12}c_{K_0,\sigma}c_{S,\sigma}\).
Iteration gives \eqref{eq:ftb-carries}.  This invocation is the actual
integer-reading theorem, including its key-unit and level-three eta
nonvanishing inputs; it is not an extension theorem for arbitrary
rational sections.

At a valued point on \(Z\), \eqref{eq:ftb-tau-bound} gives
\(v_{\mathrm{norm}}(f)\geq\phi/(625e_0)\).
The product law and the bounds for \(K_0,S\) give
\[
 v_{\mathrm{norm}}(\theta)
     \geq24\phi/(625e_0)+\rho_w
     \geq8\phi/(2500e_0)-8e_w .
\]
The last inequality uses \(\phi/e_0\geq0\).
At a point off \(Z\), the factor \(S\), supported on \(Z\), makes
\(\theta=0\), which also gives the required extended-real inequality.

The exact norm multiplication law is
\[
 \operatorname{homN}_{G_0}(\theta;y)
       =a_\sigma(y)^{12}c_{K_0,\sigma}c_{S,\sigma}.
\]
Raise the second inequality of \eqref{eq:ftb-carry-analytic} to
the eighth power.  Since
\[
 e^{16t_\sigma}
       =(e^\kappa c_{S,\sigma})^2e^{8\beta},
\]
it proves \eqref{eq:ftb-infbd}.  The source
\(\operatorname{goodOf}(\bar\sigma)=\operatorname{goodOf}(\sigma)\)
law and the constant-norm law for \(S\) imply
\(t_{\bar\sigma}=t_\sigma\).  Finally substitute
\eqref{eq:ftb-carry-choice} into the definition of cost and use
\eqref{eq:ftb-flat-carry-deg}.  The finite and infinite terms are
finite sums, so this gives \eqref{eq:ftb-carry-cost} without a
limiting argument.
\end{proof}

\begin{lemma}[The numerical cost absorption]\label{lem:ftb-carry-cost-numerical}
Suppose \(f,r,L\geq1\), \(p,q\geq0\), \(p+q=r\),
\(|\omega|\leq3L\), and \((frL)^{12}\leq P\).
Then
\[
 \frac{3r(2fr\log(2\pi))}{8}
 -\frac{12rf((p-q)\omega/2)}8\leq10P.
\]
In a common window with \(T\beta\leq C_\beta U\) and \(TP\leq U\),
the nonnegative number \(c_c=\beta/2+10P\) satisfies
\(Tc_c\leq(C_\beta+10)U\) and bounds the cost of
Proposition~\ref{prop:ftb-carry-explicit}.
\end{lemma}
\begin{proof}
The inequalities \(\log(2\pi)\leq7\) and
\(-((p-q)\omega/2)\leq3rL/2\) give the upper bound
\((42/8)fr^2+(18/8)fr^2L\).
Both \(fr^2\) and \(fr^2L\) are at most \((frL)^2\leq P\);
the displayed sum is at most \(7.5P\leq10P\).
The window assertion follows by multiplying by \(T\) and adding
the two budgets.  This uses the same \(T\) as the analytic estimate.
\end{proof}

\subsection{Spreading the reference exponents}
The source reference input has the following exact content.
There are a finite set \(S_0\) of rational primes, numbers
\(a(p)\geq0\) supported on \(S_0\), and \(s\geq0\), such that
\[
 \sum_{p\in S_0}a(p)\log p\leq s,\qquad
 \operatorname{refExp}_p(U,1/2500,G)\leq a(p)
       \quad(p\text{ prime}).                      \tag{C13}\label{eq:ftb-ref-budget}
\]
The latter supremum of normalized point depths includes, at every
algebraic integral point \(y\) over \(w\mid p\) carrying
\(\operatorname{HasPhi}(\phi)\),
\[
 \operatorname{depth}_{\mathrm{all}}(G,y)
       \leq v_w(p)a(p)+\phi/(2500e_0).              \tag{C14}\label{eq:ftb-ref-point}
\]
The depth on the left is the infimum of the normalized valuations of
the source integral cusp sections.  It uses
\(v_V(p)=e_0v_w(p)\); dividing by this equality is the definition
of the \(p\)-normalized depth.  Thus no residue-degree factor is lost.

\begin{proposition}[The same section gives the multiplier datum]
\label{prop:ftb-carry-spread}
Assume \eqref{eq:ftb-ref-budget}.  Assume also the actual valued-frame
comparison: every recurrent valuation in use admits a finite valued
extension and an algebraic integral frame carrying
\(\operatorname{HasPhi}(\phi)\), with the same normalized values
on sections.  Then
\[
                  e'_w=e_w+\sum_{p\in S_0}v_w(p)a(p)
 \tag{C15}\label{eq:ftb-spread}
\]
has finite support and
\(\operatorname{cost}(e',t)\leq\operatorname{cost}(e,t)+s\).
With \(\theta\) of Proposition~\ref{prop:ftb-carry-explicit},
the finite inequality \eqref{eq:ftb-finbd} holds for \(e'\), and
the source no-base condition holds: if any cusp section of any weight
on \(Z\) has finite recurrent valuation, some integral section of
weight \(G\) has finite valuation.  Consequently these are the
actual multiplier data of cost at most \(c_c+s\).
\end{proposition}
\begin{proof}
For every rational integer \(p>0\),
\[
 \frac1{d_K}\sum_w v_w(p)\log Nw=\log p .
\]
This is the ideal norm formula for \(p\mathcal O_K\).
It proves the cost assertion for \eqref{eq:ftb-spread}; it also
proves finite support, since an integer has only finitely many
prime ideal divisors.

Represent a recurrent valuation \(v\) by the specified extended frame.
Choose the rational prime \(p\) below \(w\).
If a cusp section on \(Z\) has finite \(v\), its nonzero evaluation
forces the point to lie on \(Z\).
For any real \(b\) strictly larger than the right-hand side of
\eqref{eq:ftb-ref-point}, the definition of infimum chooses an
integral cusp section with normalized valuation less than \(b\).
Restrict it to \(Z\); its value at the point is unchanged.
Letting \(b\) decrease gives
\[
 \operatorname{depth}_Z(G,v)\leq
 B:=v_w(p)a(p)+\phi/(2500e_0)<+\infty .
\]
This proves the no-base assertion: if every integral section had
value \(+\infty\), their infimum could not be at most \(B\).
Moreover, \eqref{eq:ftb-carry-finite} and
\(e'_w\geq e_w+v_w(p)a(p)\) imply
\[
 v(\theta)\geq8B-8e'_w
       \geq8\operatorname{depth}_Z(G,v)-8e'_w.
\]
Transport back using equality of normalized values, including
the values of \(\theta\) and of every integral section.
The infinite bound and the carrying property use the same
\(\theta,t\) and are unchanged.  If the reference budget is
\(Ts\leq C_{\mathrm{ref}}U\), addition gives total budget
\(T(c_c+s)\leq(C_\beta+10+C_{\mathrm{ref}})U\).
\end{proof}

\paragraph{Exact analytic and geometric frontiers.}
The construction above translates the carrying and spreading
calculations, with their expanded inputs.  The source produces
\eqref{eq:ftb-carry-analytic} from
\texttt{NormSqThetaConstantSpanningClassesLeExp},
\texttt{FrameExponentialVolumeLarge},
\texttt{FrameExponentialVolumeSmall}, and
\texttt{EtaBoundLengthUnit}.  In the proof the common theta parameter
is \(P_\eta=\operatorname{EtaAn.p16}(E,r,g_0g_1,m,5,1)\),
and \(\beta=P_\eta+B_{\rm small}\).
The theta bound gives \(a\leq\exp(P_\eta/625)\) everywhere and a
stronger exponent \(a\leq e^z\) in the small-first-minimum regime.
The respective Bergman lower bounds give \(G_{\rm Berg}\geq e^t\)
with \(3z\leq P_\eta+B_{\rm small}+t\).
Hence
\(a^3\leq e^{P_\eta+B_{\rm small}}G_{\rm Berg}\) in both regimes.
The length bounds for \(P_\eta\) and \(B_{\rm small}\) add with the
same window length \(T\).  These four exact producer estimates
are retained here as inputs until their proofs are translated.
The finite theta input is the same-model integral expansion,
primitive normalization and Heisenberg lift construction in
\eqref{eq:ftb-theta-coset-val}; the recurrent-frame input is
\texttt{exists_ext_hasPhi}.  The admissible shifted tuple and nonzero theta rule use
\texttt{CardLevGrpAndSpan}, fpqc descent of the theta expression, and
\texttt{ExistsShiftNullValNeZero}; the last is an actual theta-value
nonvanishing input.  The exponent-one key unit, the
integer-reading/cusp multiplication theorem, and the finite reference
budget \eqref{eq:ftb-ref-budget} are also retained with their exact
scope.  The arithmetic lemmas of this section do not by themselves
assert these geometric and analytic suppliers.
